\documentclass[10pt,a4paper]{article}
\usepackage{nexusS5}
\setnexusheader{Entrée en Troisième — Stage de pré-rentrée}{Séance 5/5 — Elyes KEFI}
\begin{document}
\nexuscover{SÉANCE 5 --- RÉINVESTIR, CONSOLIDER, PRÉPARER LA RENTRÉE}{Livret de travail}{Elyes KEFI}{Consolidation ciblée et calibration}{Entrée en Troisième --- Mathématiques --- 1\,h\,15 de travail, puis 45 minutes d'évaluation dans un document séparé}
\begin{goalbox}
\textbf{Objectif personnel de cette séance.} Séparer définitivement la procédure du produit de fractions de celle de la somme, et faire porter le signe sur la constante lors d'une réduction.
\par\medskip
\textbf{Ce livret couvre uniquement les 75 premières minutes.} L'évaluation finale est remise ensuite, dans un document distinct.
\end{goalbox}
\begin{methodbox}
\textbf{Le fil des cinq séances}
\begin{itemize}[leftmargin=13mm,labelsep=3mm,itemsep=1.5pt]
\item[\textbf{\color{Navy}S1}] Sécuriser le moteur numérique (nombres, fractions, puissances) ;
\item[\textbf{\color{Navy}S2}] Passer du calcul à l'algèbre (calcul littéral, équations) ;
\item[\textbf{\color{Navy}S3}] Construire et justifier (Pythagore, Thalès) ;
\item[\textbf{\color{Navy}S4}] Choisir un rapport trigonométrique ;
\item[\textbf{\color{Navy}S5}] aujourd'hui : réactiver, consolider deux axes personnels, faire le pont vers la Troisième, puis mesurer.
\end{itemize}
\end{methodbox}
\sectiontitle{Feuille de route des 75 minutes}
\rowcolors{2}{SoftBlue}{white}
\par\noindent\begin{tabularx}{\linewidth}{>{\bfseries}p{20mm} X X}
\rowcolor{Navy}\color{white}Temps & \color{white}Travail & \color{white}Preuve attendue \\
0--10 min & Réactiver les cinq domaines du stage et repérer immédiatement les automatismes qui se sont dégradés. & Cinq réponses écrites et le contrôle utilisé. \\
10--35 min & Produit et quotient de fractions & Trois réussites autonomes sur l'axe prioritaire. \\
35--55 min & Réduire une expression avec des constantes signées & Deux réussites autonomes sur le second axe. \\
55--70 min & De la Quatrième à la Troisième : le double produit et l'entrée dans les fonctions & Une production reliant un acquis de Quatrième à un attendu de Troisième. \\
70--75 min & Cadrage méthodologique, sans aucune information sur le contenu de l'évaluation. & Deux engagements écrits avant l'évaluation. \\
\end{tabularx}
\begin{graybox}\small\textbf{Règle de la séance.} J'écris ma réponse \emph{et} le contrôle qui la rend défendable. Une réponse sans contrôle reste une hypothèse. La ligne « Certitude » se remplit \emph{avant} toute vérification.\end{graybox}
\par\vspace{2mm}
\phasehead{Phase 1 --- Réactivation}{0--10 min}{Réactiver les cinq domaines du stage et repérer immédiatement les automatismes qui se sont dégradés.}
Sept minutes seules, sans carte de rappel, puis trois minutes de mise en commun. Écrire la réponse \emph{et} le contrôle utilisé.
\begin{enumerate}
\item Calculer $(-6)\times(-4) + (-3)\times 5$.
\shortlines{2}
\item Calculer $\dfrac{3}{4}\times\dfrac{8}{9}$ et donner le résultat simplifié.
\shortlines{2}
\item Réduire $5x - 8 + 3x + 2$, puis contrôler pour $x = 1$.
\shortlines{2}
\item $ABC$ est rectangle en $A$ avec $AB = 9$ cm et $AC = 12$ cm. Calculer $BC$ en citant le théorème.
\shortlines{2}
\item Calculer la moyenne de la série $7$ ; $11$ ; $9$ ; $13$, puis vérifier qu'elle est comprise entre le minimum et le maximum.
\shortlines{2}
\end{enumerate}
\certline
\aideline
\needspace{62mm}\par\vspace{3mm}
\phasehead{Phase 2 --- Consolidation, axe prioritaire}{10--35 min}{Produit et quotient de fractions}
\begin{methodbox}
\textbf{Deux procédures à ne pas mélanger.} Pour un \textbf{produit} : numérateurs entre eux, dénominateurs
entre eux, sans dénominateur commun. Pour un \textbf{quotient} : multiplier par l'inverse.
Pour une \textbf{somme} : réduire d'abord au même dénominateur.
\textbf{Contrôle :} estimer l'ordre de grandeur avant de calculer.
\end{methodbox}
\begin{graybox}
\textbf{Exemple guidé.} Calculer $\dfrac{3}{8}\times\dfrac{4}{9}$.
\begin{enumerate}
\item Simplifier avant de multiplier : $3$ et $9$ par $3$ ; $4$ et $8$ par $4$.
\item Il reste $\frac{1}{2}\times\frac{1}{3}=\frac{1}{6}$.
\item Contrôler : $\frac38$ vaut environ $0{,}38$ et $\frac49$ environ $0{,}44$ ; leur produit vaut environ $0{,}17$, soit $\frac16$.
\end{enumerate}
\end{graybox}
\subsectiontitle{À toi}
\begin{enumerate}
\needspace{18mm}
\item Calculer $\dfrac{5}{7}\times\dfrac{14}{15}$ et simplifier.
\worklines{2}
\needspace{18mm}
\item Calculer $\dfrac{3}{10}\times\dfrac{25}{9}$ et simplifier.
\worklines{2}
\needspace{18mm}
\item Calculer $\dfrac{8}{15}\div\dfrac{4}{5}$.
\worklines{3}
\needspace{18mm}
\item Calculer $\dfrac{9}{16}\div\dfrac{3}{8}$.
\worklines{3}
\needspace{18mm}
\item Calculer $\dfrac{2}{3}\times\dfrac{6}{5}$ puis $\dfrac{2}{3}+\dfrac{6}{5}$, et expliquer pourquoi les deux méthodes diffèrent.
\worklines{4}
\needspace{18mm}
\item Un élève écrit $\dfrac{2}{3}\times\dfrac{9}{4}=\dfrac{8}{27}$. Nommer l'erreur, corriger et donner un contrôle par estimation.
\worklines{4}
\end{enumerate}
\begin{successbox}\textbf{Critère de réussite de cette phase :} cinq produits ou quotients exacts consécutifs, dont deux comportant une simplification, avec une estimation écrite.\end{successbox}
\certline
\needspace{62mm}\par\vspace{3mm}
\phasehead{Phase 3 --- Consolidation, second axe}{35--55 min}{Réduire une expression avec des constantes signées}
\begin{methodbox}
\textbf{Le signe appartient au nombre qui le suit.} Dans $4x - 7 + 2x + 3$, les constantes sont $-7$ et $+3$,
et non $7$ et $3$. Regrouper les termes en $x$ d'un côté, les constantes de l'autre.
\textbf{Contrôle :} pour $x = 0$, l'expression réduite doit rendre exactement la somme des constantes.
\end{methodbox}
\subsectiontitle{À toi}
\begin{enumerate}
\needspace{18mm}
\item Réduire $9x - 6 + 2x - 5$.
\worklines{2}
\needspace{18mm}
\item Réduire $-5x + 12 + 8x - 3$.
\worklines{2}
\needspace{18mm}
\item Développer puis réduire $4(2x - 3) - 3x + 8$.
\worklines{3}
\needspace{18mm}
\item Développer puis réduire $-2(x - 6) + 5x - 9$.
\worklines{3}
\end{enumerate}
\begin{successbox}\textbf{Critère de réussite de cette phase :} quatre réductions consécutives exactes avec constantes signées, chacune contrôlée par $x = 0$ ou $x = -1$.\end{successbox}
\certline
\needspace{62mm}\par\vspace{3mm}
\phasehead{Phase 4 --- Réinvestissement}{55--70 min}{Découverte guidée : construire, à partir de la distributivité et des équations déjà travaillées, deux notions nouvelles de Troisième — le double produit et le vocabulaire des fonctions.}

\begin{methodbox}
\textbf{Ce qui change en Troisième.} Deux éléments sont à distinguer nettement.
\par\smallskip
\textbf{Ce qui est réinvesti} — et qui a été travaillé de la séance 1 à la séance 4 : la distributivité
$k(a+b)$, la réduction d'une expression, la résolution d'une équation du premier degré, le calcul d'une aire.
\par\smallskip
\textbf{Ce qui est nouveau} — et qui appartient au programme de Troisième : le \textbf{double produit}
$(a+b)(c+d)$, et le \textbf{statut de fonction} donné à une expression littérale, avec le vocabulaire
« image » et « antécédent ». Les deux blocs ci-dessous construisent ces deux nouveautés à partir des acquis
énumérés ci-dessus.
\end{methodbox}
\begin{warningbox}\small
\textbf{Cette phase introduit des notions de l'année à venir.} Ce que tu produis ici ne sera donc
\emph{pas} comparé à ton positionnement de début de stage : il ne s'agit pas de mesurer un progrès,
mais de préparer les premières séances de septembre.
\end{warningbox}


\subsectiontitle{A. Du rectangle au double produit — 7 min}
Un rectangle a pour longueur $x+3$ et pour largeur $x+2$. On le découpe en quatre parties.
\begin{center}
\begin{tikzpicture}[x=1cm,y=1cm]
\draw[fill=SoftBlue] (0,0) rectangle (3,2);
\draw[fill=SoftGold] (3,0) rectangle (4.5,2);
\draw[fill=SoftGray] (0,2) rectangle (3,3);
\draw[fill=green!8] (3,2) rectangle (4.5,3);
\draw[thick] (0,0) rectangle (4.5,3);
\node at (1.5,1) {$x \times x$};
\node at (3.75,1) {$3 \times x$};
\node at (1.5,2.5) {$2 \times x$};
\node at (3.75,2.5) {$3\times 2$};
\node[below] at (1.5,0) {$x$};
\node[below] at (3.75,0) {$3$};
\node[left] at (0,1) {$x$};
\node[left] at (0,2.5) {$2$};
\end{tikzpicture}
\end{center}
\begin{enumerate}
\item Écrire l'aire totale sous la forme d'un produit : \rule{55mm}{0.32pt}
\item Écrire l'aire totale comme la somme des quatre parties, puis réduire.
\shortlines{2}
\item Contrôler l'égalité obtenue pour $x = 4$ : les deux écritures doivent donner le même nombre.
\shortlines{1}
\item Développer de la même manière $(x+5)(x+1)$.
\shortlines{1}
\end{enumerate}

\subsectiontitle{B. D'un programme de calcul à une fonction — 8 min}
Programme : « choisir un nombre, le multiplier par $2$, puis retrancher $3$ ».
\begin{enumerate}
\item Écrire l'expression obtenue en partant de $x$. On la note $f(x)$. \rule{45mm}{0.32pt}
\item Calculer $f(5)$. On dit que $7$ est \rule{28mm}{0.32pt} de $5$ par $f$. \quad Calculer $f(-2)$.
\shortlines{1}
\item Quel nombre faut-il choisir pour obtenir $9$ ? Écrire l'équation, la résoudre, puis conclure :
$6$ est \rule{28mm}{0.32pt} de $9$ par $f$.
\worklines{2}
\item Compléter le tableau de valeurs.
\end{enumerate}
\par\noindent\begin{tabularx}{\linewidth}{|>{\bfseries}p{18mm}|X|X|X|X|X|}
\hline
$x$ & $-2$ & $0$ & $1$ & $3$ & $5$ \\ \hline
$f(x)$ & & & & & \\ \hline
\end{tabularx}

\begin{successbox}
\textbf{Preuve attendue :} les mots « image » et « antécédent » sont employés au bon endroit, et le calcul
d'un antécédent passe par une équation, non par un essai au hasard.
\end{successbox}

\needspace{62mm}\par\vspace{3mm}
\phasehead{Phase 5 --- Avant l'évaluation}{70--75 min}{Cadrage méthodologique. Aucun contenu de l'évaluation n'est annoncé.}

\begin{methodbox}
\textbf{Cinq règles, et rien d'autre.} Ce cadrage ne porte sur aucun contenu de l'évaluation.
\begin{enumerate}
\item \textbf{Gérer le temps.} Trois parties : environ 10 min, 20 min et 15 min. Un regard sur l'horloge à la fin de chaque partie suffit.
\item \textbf{Justifier.} Écrire la relation, la propriété ou la règle avant le calcul. Un résultat seul ne montre pas ce que tu sais.
\item \textbf{Contrôler.} Avant de passer à la suite : ordre de grandeur, substitution, unité, ou vérification par une seconde voie.
\item \textbf{Lire jusqu'au bout.} Souligner la question réellement posée et les données effectivement fournies.
\item \textbf{Passer et revenir.} Une question laissée de côté puis reprise vaut mieux qu'une réponse écrite au hasard.
\end{enumerate}
\end{methodbox}

\begin{graybox}
\textbf{La certitude.} Elle est demandée à trois moments seulement, comme au test de positionnement de début de stage :
\conf\ . Elle n'entre pas dans la note. Elle sert à distinguer ce que tu sais de ce que tu crois savoir.
\end{graybox}

\subsectiontitle{Mon engagement d'avant-évaluation}
\par\vspace{1mm}
\noindent\textbf{Le contrôle que j'écrirai systématiquement :}\ \rule{\dimexpr\linewidth-76mm\relax}{0.32pt}
\par\vspace{3mm}
\noindent\textbf{La question que je traiterai en dernier :}\ \rule{\dimexpr\linewidth-68mm\relax}{0.32pt}
\par\vspace{1mm}

\begin{warningbox}\textbf{Fin du livret de travail.} L'évaluation finale est un document séparé, remis par l'enseignant à l'issue de ces 75 minutes.\end{warningbox}
\end{document}
